%  LaTeX support: latex@mdpi.com 
%  For support, please attach all files needed for compiling as well as the log file, and specify your operating system, LaTeX version, and LaTeX editor.

%=================================================================
\documentclass[information,article,submit,pdftex,moreauthors]{Definitions/mdpi} 

%=================================================================

% MDPI internal commands
\firstpage{1} 
\makeatletter 
\setcounter{page}{\@firstpage} 
\makeatother
\pubvolume{1}
\issuenum{1}
\articlenumber{0}
\pubyear{2023}
\copyrightyear{2023}
%\externaleditor{Academic Editor: Firstname Lastname}
\datereceived{7 March 2023} 
%\daterevised{} 
\dateaccepted{} 
\datepublished{} 
\hreflink{https://doi.org/} % If needed use \linebreak
%\doinum{}

%=================================================================
% Add packages and commands here. The following packages are loaded in our class file: fontenc, inputenc, calc, indentfirst, fancyhdr, graphicx, epstopdf, lastpage, ifthen, lineno, float, amsmath, setspace, enumitem, mathpazo, booktabs, titlesec, etoolbox, tabto, xcolor, soul, multirow, microtype, tikz, totcount, changepage, attrib, upgreek, cleveref, amsthm, hyphenat, natbib, hyperref, footmisc, url, geometry, newfloat, caption

\graphicspath{{figures/}} % path to folder with figures
\usepackage{subcaption}
\usepackage{listings}
\usepackage{xcolor}
\usepackage{gensymb} % for \textdegree
\usepackage{tabularx}
\usepackage{makecell}

%=================================================================
% Full title of the paper (Capitalized)
%\Title{Land Features Recognition in Qena Bend of the Nile River, Egypt Using Unsupervised Classification of the Landsat Images by R}
\Title{Recognizing the Wadi Fluvial Structure and Stream Network in Qena Bend of the Nile River, Egypt on Landsat 8-9 OLI Images}

% MDPI internal command: Title for citation in the left column
\TitleCitation{Recognizing the Wadi Fluvial Structure and Stream Network in Qena Bend of the Nile River, Egypt on Landsat 8-9 OLI Images}

% Author Orchid ID: enter ID or remove command
\newcommand{\orcidauthorA}{0000-0002-5759-1089} % Polina Add \orcidA{} behind the author's name
\newcommand{\orcidauthorB}{0000-0002-6461-1551} % Olivier Add \orcidB{} behind the author's name

% Authors, for the paper (add full first names)
\Author{Polina Lemenkova $^{1}$*\orcidA{} and Olivier Debeir $^{1}$\orcidB{}}

%\longauthorlist{yes}

% MDPI internal command: Authors, for metadata in PDF
\AuthorNames{Polina Lemenkova and Olivier Debeir}

% MDPI internal command: Authors, for citation in the left column
\AuthorCitation{Lemenkova, P.; Debeir, O.}
% If this is a Chicago style journal: Lastname, Firstname, Firstname Lastname, and Firstname Lastname.

% Affiliations / Addresses (Add [1] after \address if there is only one affiliation.)
\address{%
$^{1}$ \quad Laboratory of Image Synthesis and Analysis (LISA), École polytechnique de Bruxelles (Brussels Faculty of Engineering), Université Libre de Bruxelles (ULB). Building L, Campus du Solbosch, ULB -- LISA CP165/57, Avenue Franklin D. Roosevelt 50, Brussels 1050, Belgium.
}

% Contact information of the corresponding author
\corres{Correspondence: polina.lemenkova@ulb.be; Tel.: +32 471 86 04 59 (P.L.)}

% Abstract
\abstract{With remote sensing data processing methods widely prevalent, the ability to handle changes in spatial data and evaluate distribution of diverse landforms across target areas in the datasets becomes important. One way to approach this problem is through satellite image processing. In this paper, we primarily focus on the methods of the unsupervised classification of the Landsat OLI/TIRS images covering the region of Qena governorate in Upper Egypt. We present a case of extracting information from the space-borne data for analysis of distribution of wadi streams over several years: 2013, 2015, 2016, 2019, 2022, and 2023. Six images were processed using computer vision methods of R. The results of k-means clustering of each scene were compared to visualise changes in landforms retrieved from the images covering Qena Bend of the Nile River. }

% Keywords
\keyword{Africa; k-means clustering; image processing; remote sensing; unsupervised classification; programming; visualization; modelling; computer vision; cartography} 

\PACS{91.10.Da; 91.10.Jf; 91.10.Sp; 91.10.Xa; 96.25.Vt; 91.10.Fc; 95.40.+s; 95.75.Qr; 95.75.Rs; 42.68.Wt}
\MSC{86A30; 86-08; 86A99; 86A04}
\JEL{Y91; Q20; Q24; Q23; Q3; Q01; R11; O44; O13; Q5; Q51; Q55; N57; C6; C61}

%%%%%%%%%%%%%%%%%%%%%%%%%%%%%%%%%%%%%%%%%%
\begin{document}

\section{Introduction}

Satellite image classification methods are frequently used to produce maps of land cover types and recognise land features. Example  of such applications include maps of mineral resources \cite{ABDELKAREEM2018682,1294518,rs11121450}, environmental assessment \cite{615837,Lemenkova202227}, land cover/land use classification \cite{516791,Lemenkova202229}, geologic analysis \cite{6723607,KAMEL2022104420} and urban mapping \cite{MOHAMED2019103507}. The success of the remote sensing in cartographic data processing and mapping is based on the effectiveness of the satellite imagery as a source of geoinformation. Thus, using high-resolution space-borne data improves the quality of spatial visualization \cite{ABDALLA20128,Lemenkova202230}, provides precious source of information in remotely located places and inaccessible areas such as deserts \cite{KAMEL2022104420,6350376}. Moreover, comparison of several images on target region collected in different time spans enables time series analysis aimed at detecting environmental changes \cite{9171725,SINGH2020423,Hashim}. 

Applications of the remote sensing data in Earth sciences include monitoring the desertification and land degradation \cite{Badreldin,Kawy}, analysis of deforestation \cite{516791}, evaluation of salinization \cite{8949893} and soil erosion \cite{Gad2020}, hydrological modelling and quantification of the flooded areas \cite{TAHA2017157,ELSADEK2019377} and many more. Spectral patterns of the multispectral Landsat images are especially widely used in geosciences as a source of spatial information for land cover mapping and recognition of land features \cite{8486245}. Natural resource monitoring 

\pagebreak
\section{Prior Work}

The region of Qena Bend, situated in Upper Egypt (Figure \ref{fig01}), is affected both by the hydrology of the Nile River and by the influence of the Eastern Desert. A special feature of the climate of Eastern Desert includes the occasional extreme flash floods -- one of the most severe natural disasters caused by the torrential rainfalls and resulting in the disastrous surface runoff \cite{Moawad}. At the same time, the geomorphology of this area is notable for a complex drainage network of wadi aquifer systems \cite{Alkholy,AbdelMoneim}. This presents an interrelated complex of the ephemeral rivers dry during most of the year and containing water during flash floods and heavy showers and rains during spring and autumn annually \cite{ElFakharany}. The cumulative effects from wadi network and arid and semi-arid climate of the Eastern Desert region results in occasional flash flood events and irregular sand dust storms with different durations and intensities.

\begin{figure}[H]
\centering
	\includegraphics[width=14.0 cm]{Fig_01.jpg}
	\caption{Topographic map of Egypt. Target study area of Qena Bend region is shown in red rotated square. Software: GMT v. 6.1.1. Data source: GEBCO/SRTM. Cartography source: authors.
	\label{fig01}}
\end{figure}

Flash floods occurs in diverse regions across in Egypt including Upper Egypt with Eastern Desert \cite{ElHaddad,Elsadek}, southern Red Sea Coast \cite{Abdalla2014}, Sinai Peninsula \cite{ELOSTA2021101522} and Gulf of Suez \cite{AbdelLattif}. The social consequences of flash floods include damaged infrastructure, transport ways and buildings, and occasional victims \cite{Moawad}. El-Magd et al. report the catastrophic events of flash floods in Qena governorate recorded in three consecutive events in 2014--2016 \cite{ElMagd}. 
The occurrence of flood hazards in Upper Egypt point at the significance of the application of geoinformation for visualization and mapping the regions prone to flood risk. The remote sensing data and advanced methods of spatial data processing present robust approaches to evaluation of relief and extracting information on relief and topographic characteristics from the satellite images using methods of computer vision. Such approaches can be used, for instance for vulnerability mapping or identification of areas with high geomorphological hazards and expose to floods \cite{nhess-9-751-2009}.

The problem of morphometric modelling and visualization was investigated in cartography, geoinformatics and Earth science communities. Hamdan. and Khozyem \cite{Hamdan} applied ASTER data and GIS techniques to delineate drainage networks in Wadi El-Mathula watershed in Central Eastern Desert. The geometry and morphology of relief, the characteristics of topographic parameters such as slope, aspect and hillshade can be extracted as morphometric characteristics for analysis of watershed using 3D mapping \cite{Lemenkova202301}. Such information can be used for complex hydrological analysis, for instance, predicting wadi flash floods using morphometric parameters that represent the topographic and drainage characteristics of the basins using geographic information system (GIS) \cite{w9070553}. Basin and drainage network,

Extracting information from the satellite images for analysis of drainage wadi networks enables to define basin boundaries and stream channel density, perform visualization of slope and flow directions for modelling of the hydrological processes in the past to analyse geologic setting in present \cite{abdelkareem2013integration,AbdelkareemFarouk}. Further, the essential geographic information collected from the satellite imagery is useful for operative flash flood monitoring through integration of the data on physiographic features. Other applications of the remote sensing data include modelling lithological, geological, and structural features \cite{BESHR2021999}, mapping mineral resources \cite{Kamal} and detecting wadi channels based on the combination of the topographic maps and satellite images \cite{Moubark}.

%%%%%%%%%%%%%%%%%%%%%%%%%%%%%%%%%%%%%%%%%%
\newpage
\section{Materials and Methods}

\subsection{Data}

Multispectral satellite images such as Landsat TM/ETM+, and OLI/TIRS) provide a valuable source of remote sensing data widely used in environmental studies of Egypt \cite{Elfadaly,Sultanetal1986,Sultanetal1987}. In this work, we use the two major types of data Figure \ref{fig02}: the topographic grid for mapping the study area and a series of six satellite images presented by the Landsat 8-9 OLI/TIRS scenes. The high-resolution topographic data were used for inspection of the relief in the region, while the remote sensing data were used for comparative analysis of changes in the wadi channels in the Qena Bend area and surroundings by image processing and k-means clustering algorithms using R. 

\begin{figure}[H]
\centering
	\includegraphics[width=14.0 cm]{Fig_02}
	\caption{Data sources: summary of the materials used in this study. Software: R version 4.2.2, 'Mermaid' library. Flowchart source: authors.
	\label{fig02}}
\end{figure}

The detailed information on the satellite images is summarised in Table \ref{tab1}. 

\begin{table}[H]
\small
\caption{Satellite images used for computing the VIs: Landsat-8 OLI/TIRS collected from the USGS. \textsuperscript{1}.\label{tab1}}
	\begin{adjustwidth}{-\extralength}{0cm}
%		\newcolumntype{C}{>{\centering\arraybackslash}X}
		\begin{tabularx}{\fulllength}{ p{1.6cm} | p{1.8cm} | p{7.0cm} | p{4.0cm} | p{2.0cm} }
			\toprule
			\textbf{Date} & \textbf{Spacecraft} & \textbf{Landsat Product ID} & \textbf{Scene ID} & \textbf{Cloudiness} \\
			\midrule
			2013/11/16 & Landsat 8 & \makecell{LC08\_L2SP\_175042\_20131116\_20200912\_02\_T1} & LC81750422013320LGN01 & 0.89 \\
			2015/11/22 & Landsat 8 & \makecell{LC08\_L2SP\_175042\_20151122\_20200908\_02\_T1} & LC81750422015326LGN01 & 0.05 \\
			2016/11/08 & Landsat 8 & \makecell{LC08\_L2SP\_175042\_20161108\_20200905\_02\_T1} & LC81750422016313LGN01 & 0.31 \\
			2019/11/17 & Landsat 8 & \makecell{LC08\_L2SP\_175042\_20191117\_20200825\_02\_T1} & LC81750422019321LGN00 & 1.16 \\
			2022/11/09 & Landsat 8 & \makecell{LC08\_L2SP\_175042\_20221109\_20221121\_02\_T1} & LC81750422022313LGN00 & 0.02 \\
			2023/01/20 & Landsat 9 & \makecell{LC09\_L2SP\_175042\_20230120\_20230122\_02\_T1} & LC91750422023020LGN01 & 0.32 \\	
			\bottomrule
		\end{tabularx}
	\end{adjustwidth}
	\noindent{\footnotesize{\textsuperscript{1} The Sensor ID is common for all the scenes: Landsat 8-9 OLI/TIRS (Operational Land Imager and Thermal Infrared Sensor), Collection 2 Level-2. Path/row parameters common for all the images: 175/42. Coverage: Qena Bend, the Nile River, Upper Egypt. Image source: the USGS \cite{EROS}}}
\end{table}

\begin{figure}[H]
\makebox[1.00\linewidth][r]{%
	\begin{subfigure}[b]{.40\textwidth}
		\centering
			\includegraphics[width=0.87\textwidth]{Fig_03a.jpg}
		\caption{2013}
	\end{subfigure}%
	\begin{subfigure}[b]{.40\textwidth}
		\centering
		\includegraphics[width=0.87\textwidth]{Fig_03b.jpg}
		\caption{2015}
	\end{subfigure}%
	\begin{subfigure}[b]{.40\textwidth}
		\centering
		\includegraphics[width=0.87\textwidth]{Fig_03c.jpg}
		\caption{2018}
	\end{subfigure}%
}\\
\vfill \vspace{1mm}
\makebox[1.00\linewidth][r]{%
	\begin{subfigure}[b]{.40\textwidth}
		\centering
			\includegraphics[width=0.87\textwidth]{Fig_03d.jpg}
		\caption{2020}
	\end{subfigure}%
	\begin{subfigure}[b]{.40\textwidth}
		\centering
		\includegraphics[width=0.87\textwidth]{Fig_03e.jpg}
		\caption{2021}
	\end{subfigure}%
	\begin{subfigure}[b]{.40\textwidth}
		\centering
		\includegraphics[width=0.87\textwidth]{Fig_03f.jpg}
		\caption{2022}
	\end{subfigure}%
}\caption{Qena Bend of the Nile River, Egypt visible on Landsat 8-9 OLI/TIRS images in natural RGB colour for six years: (\textbf{a}) 2013/12/20, (\textbf{b}) 2015/12/26, (\textbf{c}) 2018/02/21, (\textbf{d}) 2020/01/22, (\textbf{e}) 2021/12/18, (\textbf{f}) 2022/12/29.}\label{fig:03}
\end{figure}

%----------- subsection ----------->
\subsection{Methods}

\begin{figure}[H]
\centering
	\includegraphics[width=12.0 cm]{Fig_04}
	\caption{Methodological workflow of the processes and research steps applied in this study. Software: R version 4.2.2, 'DiagrammeR' library. Flowchart source: authors.
	\label{fig04}}
\end{figure}

%%%%%%%%%%%%%%%%%%%%%%%%%%%%%%%%%%%%%%%%%%
\pagebreak
\section{Results}

\begin{table}[H]
\small
\caption{Results of the k-means classification of the Landsat-8 OLI images for Qena Bend area\textsuperscript{1}.\label{tab2}}
	\begin{adjustwidth}{-\extralength}{0cm}
		\begin{tabularx}{\fulllength}{ C | C | C | C | C | C | C }
			\toprule
			\textbf{Class ID} & \textbf{2013} & \textbf{2015} & \textbf{2016} & \textbf{2019} & \textbf{2022} & \textbf{2023} \\
			\midrule
			1 & 6237758.6 & 25058427 & 11830859 & 11050520 & 33773365.2 & 22414794 \\
			2 & 26530972.6 & 3829403 & 6847520 & 7687698 & 21467288.5 & 4080038 \\
			3 & 46356528.0 & 23160871 & 13056437 & 31925587 & 466472.5 & 13918989 \\
			4 & 10314256.8 & 7458695 & 12594525 & 3359878 & 3547592.0 & 1863680 \\
			5 & 7861914.8 & 9909609 & 3012708 & 11101016 & 7714730.1 & 6083803 \\
			6 & 625035.3 & 4568955 & 2195275 & 14290386 & 9346294.4 & 14743313 \\
			7 & 22694195.3 & 3019950 & 12277811 & 14898378 & 6978872.8 & 6370890 \\
			8 & 5804917.1 & 1942973 & 10144633 & 5386949 & 16906933.2 & 6053497 \\
			9 & 7320502.4 & 13450745 & 3220846 & 5501193 & 8429786.2 & 11952344 \\
			10 & 6121676.5 & 10228437 & 4195904 & 10594134 & 11161798.4 & 10517903 \\	
			\bottomrule
		\end{tabularx}
	\end{adjustwidth}
	\noindent{\footnotesize{\textsuperscript{1} Sum of squares by clusters within computed cluster centroids.}}
\end{table}

\begin{figure}[H]
\makebox[1.00\linewidth][r]{%
	\begin{subfigure}[b]{.40\textwidth}
		\centering
			\includegraphics[width=0.87\textwidth]{Fig_04a}
		\caption{2013: Bands 432}
	\end{subfigure}%
	\begin{subfigure}[b]{.40\textwidth}
		\centering
		\includegraphics[width=0.87\textwidth]{Fig_04b}
		\caption{2013: Bands 543}
	\end{subfigure}%
	\begin{subfigure}[b]{.40\textwidth}
		\centering
		\includegraphics[width=1.00\textwidth]{Fig_04c}
		\caption{2013: Clustering}
	\end{subfigure}%
}\\
\vfill \vspace{1mm}
\makebox[1.00\linewidth][r]{%
	\begin{subfigure}[b]{.40\textwidth}
		\centering
			\includegraphics[width=0.87\textwidth]{Fig_04d}
		\caption{2015: Bands 571}
	\end{subfigure}%
	\begin{subfigure}[b]{.40\textwidth}
		\centering
		\includegraphics[width=0.87\textwidth]{Fig_04e}
		\caption{2015: Bands 632}
	\end{subfigure}%
	\begin{subfigure}[b]{.40\textwidth}
		\centering
		\includegraphics[width=0.92\textwidth]{Fig_04f}
		\caption{2015: Clustering}
	\end{subfigure}%
}\caption{Clustering of the Landsat 8-9 OLI/TIRS images Qena Bend of the Nile River, Egypt visible on in natural RGB colour for six years: (\textbf{a}) 2013 RGB in 432 Bands, (\textbf{b}) 2013 RGB in 543 Bands, (\textbf{c}) 2013 unsupervised classification, (\textbf{d}) 2015 RGB in 571 Bands, (\textbf{e}) 2015 RGB in 632 Bands, (\textbf{f}) 2015 RGB: Clustering.}\label{fig:04}
\end{figure}

\begin{table}[H]
\small
\caption{Pairwise comparison of the k-means classification for Bands 5-4-3 of the Landsat-8 OLI/TIRS images for 2013 and 2015 in Qena Bend region. \textsuperscript{1}.\label{tab2}}
	\begin{adjustwidth}{-\extralength}{0cm}
		\begin{tabularx}{\fulllength}{ C || C | C | C || C | C | C }
			\toprule
			\textbf{Class ID} & \textbf{\makecell{2013 B5}} & \textbf{2013 B4} & \textbf{2013 B3} & \textbf{2015 B5} & \textbf{2015 B4} & \textbf{2015 B3} \\
			\midrule
			1 & 21019.08 & 18348.92 & 15722.58 & 19819.50 & 17226.143 & 14765.14 \\
			2 & 21956.38 & 19148.12 & 16294.50 & 12087.00 & 10917.333 & 10115.00 \\
			3 & 26892.38 & 23105.00 & 19120.92 & 18736.17 & 9956.333 & 10142.17 \\
			4 & 14308.50 & 13163.75 & 11730.75 & 24835.46 & 21496.385 & 17990.38 \\
			5 & 19246.50 & 16942.69 & 14626.81 & 23315.56 & 20178.688 & 17043.31 \\
			6 & 17656.18 & 15405.09 & 13460.00 & 26302.78 & 22807.222 & 19156.67 \\
			7 & 23465.45 & 20280.18 & 17179.82 & 15028.50 & 13507.500 & 12109.33 \\
			8 & 16549.67 & 10211.17 & 10132.83 & 29412.00 & 25219.500 & 20731.50 \\
			9 & 22484.00 & 19481.25 & 16802.50 & 17737.46 & 15833.385 & 13783.23 \\
			10 & 24895.53 & 21740.20 & 18326.47 & 17737.46 & 18747.389 & 15990.33 \\	
			\bottomrule
		\end{tabularx}
	\end{adjustwidth}
	\noindent{\footnotesize{\textsuperscript{1} Clusters are computed for each band in the composite triplets: Band 5 (NIR), Band 4 (Red) and Band 3 (Green)}}
\end{table}

\pagebreak

\begin{figure}[H]
\makebox[1.00\linewidth][r]{%
	\begin{subfigure}[b]{.40\textwidth}
		\centering
			\includegraphics[width=0.87\textwidth]{Fig_05a}
		\caption{2016: Bands 564}
	\end{subfigure}%
	\begin{subfigure}[b]{.40\textwidth}
		\centering
		\includegraphics[width=0.87\textwidth]{Fig_05b}
		\caption{2016: Bands 754}
	\end{subfigure}%
	\begin{subfigure}[b]{.40\textwidth}
		\centering
		\includegraphics[width=0.90\textwidth]{Fig_05c}
		\caption{2016: Clustering}
	\end{subfigure}%
}\\
\vfill \vspace{1mm}
\makebox[1.00\linewidth][r]{%
	\begin{subfigure}[b]{.40\textwidth}
		\centering
			\includegraphics[width=0.87\textwidth]{Fig_05d}
		\caption{2019: Bands 764}
	\end{subfigure}%
	\begin{subfigure}[b]{.40\textwidth}
		\centering
		\includegraphics[width=0.87\textwidth]{Fig_05e}
		\caption{2019: Bands 753}
	\end{subfigure}%
	\begin{subfigure}[b]{.40\textwidth}
		\centering
		\includegraphics[width=0.92\textwidth]{Fig_05f}
		\caption{2019: Clustering}
	\end{subfigure}%
}\caption{Clustering of the Landsat 8-9 OLI/TIRS images Qena Bend of the Nile River, Egypt visible on in natural RGB colour for six years: (\textbf{a}) 2016 RGB in 564 Bands, (\textbf{b}) 2016 RGB in 754 Bands, (\textbf{c}) 2016 unsupervised classification, (\textbf{d}) 2019 RGB in 764 Bands, (\textbf{e}) 2019 RGB in 753 Bands, (\textbf{f}) 2019 RGB: Clustering.}\label{fig:05}
\end{figure}

\begin{table}[H]
\small
\caption{Pairwise comparison of the k-means classification for Bands 5-4-3 of the Landsat-8 OLI/TIRS images for 2016 and 2019 in Qena Bend region. \textsuperscript{1}.\label{tab2}}
	\begin{adjustwidth}{-\extralength}{0cm}
		\begin{tabularx}{\fulllength}{ C || C | C | C || C | C | C }
			\toprule
			\textbf{Class ID} & \textbf{2016 B5} & \textbf{2016 B4} & \textbf{2016 B3} & \textbf{2019 B5} & \textbf{2019 B4} & \textbf{2019 B3} \\
			\midrule
			1 & 25067.53 & 21805.82 & 18283.65 & 20721.42 & 18438.42 & 15725.33 \\
			2 & 19837.27 & 17506.45 & 15060.45 & 24198.84 & 21005.26 & 17763.84 \\
			3 & 27148.44 & 23483.89 & 19596.67 & 25536.93 & 22109.73 & 18481.20 \\
			4 & 16049.43 & 13237.00 & 11880.43 & 18257.00 & 9376.75 & 9707.50 \\
			5 & 19541.25 & 9323.50 & 9742.75 & 14476.67 & 12599.83 & 11411.17 \\
			6 & 12780.50 & 11767.00 & 10670.50 & 16992.00 & 15236.88 & 13304.75 \\
			7 & 21370.47 & 18891.65 & 16208.94 & 18829.53 & 16928.47 & 14654.00 \\
			8 & 23122.89 & 20069.28 & 16927.33 & 31113.00 & 26355.00 & 21077.00 \\
			9 & 18240.88 & 16536.62 & 14205.25 & 27920.00 & 23667.00 & 19115.00 \\
			10 & 17269.14 & 15186.43 & 13284.43 & 22736.25 & 19858.25 & 16938.06 \\	
			\bottomrule
		\end{tabularx}
	\end{adjustwidth}
	\noindent{\footnotesize{\textsuperscript{1} Clusters are computed for each band in the composite triplets: Band 5 (NIR), Band 4 (Red) and Band 3 (Green)}}
\end{table}

\pagebreak

\begin{figure}[H]
\makebox[1.00\linewidth][r]{%
	\begin{subfigure}[b]{.40\textwidth}
		\centering
			\includegraphics[width=0.87\textwidth]{Fig_06a}
		\caption{2022: Bands 631}
	\end{subfigure}%
	\begin{subfigure}[b]{.40\textwidth}
		\centering
		\includegraphics[width=0.87\textwidth]{Fig_06b}
		\caption{2022: Bands 543}
	\end{subfigure}%
	\begin{subfigure}[b]{.40\textwidth}
		\centering
		\includegraphics[width=0.90\textwidth]{Fig_06c}
		\caption{2022: Clustering}
	\end{subfigure}%
}\\
\vfill \vspace{1mm}
\makebox[1.00\linewidth][r]{%
	\begin{subfigure}[b]{.40\textwidth}
		\centering
			\includegraphics[width=0.87\textwidth]{Fig_06d}
		\caption{2023: Bands 562}
	\end{subfigure}%
	\begin{subfigure}[b]{.40\textwidth}
		\centering
		\includegraphics[width=0.87\textwidth]{Fig_06e}
		\caption{2023: Bands 753}
	\end{subfigure}%
	\begin{subfigure}[b]{.40\textwidth}
		\centering
		\includegraphics[width=0.92\textwidth]{Fig_06f}
		\caption{2023: Clustering}
	\end{subfigure}%
}\caption{Clustering of the Landsat 8-9 OLI/TIRS images Qena Bend of the Nile River, Egypt visible on in natural RGB colour for six years: (\textbf{a}) 2022 RGB in 631 Bands, (\textbf{b}) 2022 RGB in 543 Bands, (\textbf{c}) 2022 unsupervised classification, (\textbf{d}) 2023 RGB in 562 Bands, (\textbf{e}) 2023 RGB in 753 Bands, (\textbf{f}) 2023 RGB: Clustering.}\label{fig:06}
\end{figure}

\begin{table}[H]
\small
\caption{Pairwise comparison of the k-means classification for Bands 5-4-3 of the Landsat-8 OLI/TIRS images for 2022 and 2023 in Qena Bend region. \textsuperscript{1}.\label{tab2}}
	\begin{adjustwidth}{-\extralength}{0cm}
		\begin{tabularx}{\fulllength}{ C | C | C | C | C | C | C }
			\toprule
			\textbf{Class ID} & \textbf{2022 B5} & \textbf{2022 B4} & \textbf{2022 B3} & \textbf{2023 B5} & \textbf{2023 B4} & \textbf{2023 B3} \\
			\midrule
			1 & 24428.50 & 21422.20 & 18051.50 & 21513.50 & 10245.750 & 10483.000 \\
			2 & 19238.75 & 17313.00 & 14911.25 & 25528.40 & 22234.800 & 18665.200 \\
			3 & 27402.67 & 23631.83 & 19472.50 & 17918.11 & 9096.444 & 9575.333 \\
			4 & 22303.50 & 19676.83 & 16793.33 & 24549.75 & 21362.375 & 17997.875 \\
			5 & 17517.33 & 11594.67 & 10989.33 & 15965.40 & 14835.200 & 13574.400 \\
			6 & 12943.67 & 11050.00 & 10303.00 & 22865.67 & 19456.933 & 16537.733 \\
			7 & 17225.22 & 15122.22 & 13253.00 & 23488.23 & 20695.000 & 17726.000 \\
			8 & 25640.82 & 22323.82 & 18730.73 & 18460.75 & 16531.833 & 14379.417 \\
			9 & 20699.44 & 18365.81 & 15840.12 & 20155.57 & 18063.786 & 15478.143 \\
			10 & 23487.22 & 20428.72 & 17181.22 & 27425.50 & 23626.900 & 19359.200 \\	
			\bottomrule
		\end{tabularx}
	\end{adjustwidth}
	\noindent{\footnotesize{\textsuperscript{1} Clusters are computed for each band in the composite triplets: Band 5 (NIR), Band 4 (Red) and Band 3 (Green)}}
\end{table}

%%%%%%%%%%%%%%%%%%%%%%%%%%%%%%%%%%%%%%%%%%
\pagebreak
\section{Discussion}

%%%%%%%%%%%%%%%%%%%%%%%%%%%%%%%%%%%%%%%%%%
\section{Conclusions}
We have approached the problem of unsupervised classification of the satellite images by R for extracting information on the land cover types in Qena Bend region of the Upper Egypt. The application of R libraries 'raster', 'terra' and 'RStoolbox' enabled to retrieve information form the remote sensing data and analyse changes in the complex of the ephemeral riverbeds in Qena district by methods of computer vision and applied programming.

%%%%%%%%%%%%%%%%%%%%%%%%%%%%%%%%%%%%%%%%%%
\authorcontributions{Supervision, conceptualization, methodology, software, resources, funding acquisition, and project administration, O.D.; writing---original draft preparation, methodology, software, data curation, visualization, formal analysis, validation, writing---review and editing, and investigation, P.L. All authors have read and agreed to the published version of the manuscript.}

\funding{The publication was funded by the Editorial Office of \emph{Information}, MDPI, by providing 100\% discount for the APC of this manuscript. This project was supported by the Federal Public Planning Service Science Policy or Belgian Federal Science Policy Office - BELSPO (B2/202/P2/SEISMOSTORM).}

\institutionalreview{Not applicable.}

\informedconsent{Not applicable.}

\dataavailability{Not applicable} 

\acknowledgments{The authors thank the reviewers for reading and review of this manuscript.}

\conflictsofinterest{The authors declare no conflict of interest.} 

\abbreviations{Abbreviations}{
The following abbreviations are used in this manuscript:\\

\noindent 
\begin{tabular}{@{}ll}
MDPI & Multidisciplinary Digital Publishing Institute\\
DOAJ & Directory of open access journals\\
LD & Linear dichroism
\end{tabular}
}

\begin{adjustwidth}{-\extralength}{0cm}
%\printendnotes[custom] % Un-comment to print a list of endnotes

\bibliography{Qena.bib}
\reftitle{References}
\nocite{*}

%=====================================
% References, variant A: external bibliography
%=====================================


%%%%%%%%%%%%%%%%%%%%%%%%%%%%%%%%%%%%%%%%%%
\end{adjustwidth}
\end{document}
